% ============================================================
\chapter{CI/CD for Machine Learning}
\label{ch:cicd}
\index{CI/CD}
% ============================================================

\section{Why CI/CD for ML?}
\label{sec:why-cicd}

Traditional CI/CD pipelines test code and deploy applications.
ML pipelines must additionally validate data, test model quality, and
manage artifacts that are much larger than compiled binaries.

\begin{definition}[CI/CD for ML]
\textbf{Continuous Integration}\index{continuous integration} for ML
verifies that code changes do not break the training pipeline, that
data schemas are respected, and that model performance meets minimum
thresholds. \textbf{Continuous Deployment}\index{continuous deployment}
automates the promotion of validated models to staging and production
environments.
\end{definition}

\begin{center}
\begin{tikzpicture}[
  box/.style={rectangle, draw, rounded corners, minimum width=2.2cm,
    minimum height=1cm, align=center, font=\scriptsize\bfseries},
  arr/.style={-{Stealth[length=2mm]}, thick, steelblue}
]
  \node[box, fill=blue!10] (push) {Git Push};
  \node[box, fill=cyan!10, right=1.2cm of push] (lint) {Lint \&\\Format};
  \node[box, fill=cyan!10, right=1.2cm of lint] (test) {Unit\\Tests};
  \node[box, fill=green!10, right=1.2cm of test] (train) {Train\\(small)};
  \node[box, fill=green!10, right=1.2cm of train] (valid) {Model\\Validation};
  \node[box, fill=orange!10, right=1.2cm of valid] (deploy) {Deploy\\Staging};
  \node[box, fill=red!10, right=1.2cm of deploy] (prod) {Deploy\\Production};

  \draw[arr] (push) -- (lint);
  \draw[arr] (lint) -- (test);
  \draw[arr] (test) -- (train);
  \draw[arr] (train) -- (valid);
  \draw[arr] (valid) -- (deploy);
  \draw[arr] (deploy) -- (prod);

  \node[below=0.5cm of lint, font=\scriptsize, text=gray] {CI};
  \node[below=0.5cm of deploy, font=\scriptsize, text=gray] {CD};

  \draw[dashed, gray] ([xshift=0.7cm, yshift=-0.3cm]test.south east)
    -- ([xshift=0.7cm, yshift=0.8cm]test.north east);
\end{tikzpicture}
\end{center}

\begin{remark}
The key difference between CI/CD for software and CI/CD for ML is the
\textbf{model validation} step: the pipeline must verify that the
newly trained model meets performance thresholds before deployment. A
model that passes all unit tests can still be useless if accuracy has
regressed.
\end{remark}

\section{Pre-commit Hooks}
\label{sec:pre-commit}
\index{pre-commit}

Pre-commit hooks run checks \emph{before} code enters the repository,
catching issues at the earliest possible stage.

\begin{codeblock}[title=\textbf{.pre-commit-config.yaml}]
\begin{minted}{yaml}
repos:
  - repo: https://github.com/pre-commit/pre-commit-hooks
    rev: v4.5.0
    hooks:
      - id: trailing-whitespace
      - id: end-of-file-fixer
      - id: check-yaml
      - id: check-added-large-files
        args: ['--maxkb=1000']

  - repo: https://github.com/psf/black
    rev: 24.2.0
    hooks:
      - id: black

  - repo: https://github.com/astral-sh/ruff-pre-commit
    rev: v0.3.0
    hooks:
      - id: ruff
        args: [--fix]

  - repo: https://github.com/pre-commit/mirrors-mypy
    rev: v1.8.0
    hooks:
      - id: mypy
        additional_dependencies: [numpy, pandas-stubs]
\end{minted}
\end{codeblock}

\begin{codeblock}[title=\textbf{Setting up pre-commit}]
\begin{minted}{bash}
# Install pre-commit
pip install pre-commit

# Install hooks in the repository
pre-commit install

# Run on all files (first time)
pre-commit run --all-files

# Update hooks to latest versions
pre-commit autoupdate
\end{minted}
\end{codeblock}

\section{GitHub Actions for ML}
\label{sec:github-actions}
\index{GitHub Actions}

GitHub Actions is the most popular CI/CD platform for open-source ML
projects. Workflows are defined in YAML files under
\file{.github/workflows/}.

\begin{codeblock}[title=\textbf{.github/workflows/ml-ci.yml --- Complete ML CI pipeline}]
\begin{minted}{yaml}
name: ML CI Pipeline

on:
  push:
    branches: [main, develop]
  pull_request:
    branches: [main]

env:
  PYTHON_VERSION: "3.11"
  MLFLOW_TRACKING_URI: "sqlite:///mlflow.db"

jobs:
  lint-and-test:
    runs-on: ubuntu-latest
    steps:
      - uses: actions/checkout@v4

      - name: Set up Python
        uses: actions/setup-python@v5
        with:
          python-version: ${{ env.PYTHON_VERSION }}
          cache: pip

      - name: Install dependencies
        run: |
          pip install -r requirements.txt
          pip install -r requirements-dev.txt

      - name: Lint with ruff
        run: ruff check src/ tests/

      - name: Format check with black
        run: black --check src/ tests/

      - name: Type check with mypy
        run: mypy src/ --ignore-missing-imports

      - name: Unit tests
        run: pytest tests/unit/ -v --tb=short

  train-and-validate:
    needs: lint-and-test
    runs-on: ubuntu-latest
    steps:
      - uses: actions/checkout@v4

      - name: Set up Python
        uses: actions/setup-python@v5
        with:
          python-version: ${{ env.PYTHON_VERSION }}
          cache: pip

      - name: Install dependencies
        run: pip install -r requirements.txt

      - name: Set up DVC
        uses: iterative/setup-dvc@v2

      - name: Pull data with DVC
        run: dvc pull
        env:
          AWS_ACCESS_KEY_ID: ${{ secrets.AWS_ACCESS_KEY_ID }}
          AWS_SECRET_ACCESS_KEY: ${{ secrets.AWS_SECRET_KEY }}

      - name: Train model (small dataset)
        run: |
          python src/train.py \
            --config configs/ci_config.yaml \
            --max-samples 1000

      - name: Validate model performance
        run: |
          python src/evaluate.py \
            --model models/model.joblib \
            --min-accuracy 0.85 \
            --min-f1 0.80

      - name: Upload model artifact
        uses: actions/upload-artifact@v4
        with:
          name: trained-model
          path: models/model.joblib
          retention-days: 7

  integration-test:
    needs: train-and-validate
    runs-on: ubuntu-latest
    services:
      postgres:
        image: postgres:16-alpine
        env:
          POSTGRES_DB: test_db
          POSTGRES_USER: test_user
          POSTGRES_PASSWORD: test_pass
        ports:
          - 5432:5432
    steps:
      - uses: actions/checkout@v4

      - name: Set up Python
        uses: actions/setup-python@v5
        with:
          python-version: ${{ env.PYTHON_VERSION }}

      - name: Download model artifact
        uses: actions/download-artifact@v4
        with:
          name: trained-model
          path: models/

      - name: Install dependencies
        run: pip install -r requirements.txt

      - name: Run API integration tests
        run: pytest tests/integration/ -v
\end{minted}
\end{codeblock}

\section{Model Validation in CI}
\label{sec:model-validation}
\index{model validation}

Model validation ensures that newly trained models meet minimum quality
standards before deployment.

\begin{codeblock}[title=\textbf{src/evaluate.py --- Model validation script}]
\begin{minted}{python}
"""Model validation for CI pipeline."""
import argparse
import json
import sys

import joblib
import pandas as pd
from sklearn.metrics import (
    accuracy_score, f1_score, precision_score, recall_score
)


def validate_model(
    model_path: str,
    test_data_path: str = "data/test.csv",
    min_accuracy: float = 0.85,
    min_f1: float = 0.80,
) -> bool:
    """Validate model against minimum thresholds."""
    model = joblib.load(model_path)
    df = pd.read_csv(test_data_path)
    X_test = df.drop("target", axis=1)
    y_test = df["target"]

    y_pred = model.predict(X_test)

    metrics = {
        "accuracy": accuracy_score(y_test, y_pred),
        "f1": f1_score(y_test, y_pred, average="weighted"),
        "precision": precision_score(y_test, y_pred,
                                     average="weighted"),
        "recall": recall_score(y_test, y_pred,
                               average="weighted"),
    }

    # Save metrics report
    with open("metrics.json", "w") as f:
        json.dump(metrics, f, indent=2)

    print("=== Model Validation Report ===")
    for name, value in metrics.items():
        status = "PASS" if value >= locals().get(
            f"min_{name}", 0.0) else "FAIL"
        print(f"  {name}: {value:.4f} [{status}]")

    passed = (
        metrics["accuracy"] >= min_accuracy
        and metrics["f1"] >= min_f1
    )
    print(f"\nOverall: {'PASSED' if passed else 'FAILED'}")
    return passed


if __name__ == "__main__":
    parser = argparse.ArgumentParser()
    parser.add_argument("--model", required=True)
    parser.add_argument("--test-data", default="data/test.csv")
    parser.add_argument("--min-accuracy", type=float, default=0.85)
    parser.add_argument("--min-f1", type=float, default=0.80)
    args = parser.parse_args()

    success = validate_model(
        args.model, args.test_data,
        args.min_accuracy, args.min_f1,
    )
    sys.exit(0 if success else 1)
\end{minted}
\end{codeblock}

\section{DVC in CI}
\label{sec:dvc-ci}
\index{DVC!CI}

Data Version Control (DVC) integrates with CI pipelines to pull
versioned data and reproduce training pipelines.

\begin{codeblock}[title=\textbf{DVC integration in GitHub Actions}]
\begin{minted}{yaml}
- name: Set up DVC
  uses: iterative/setup-dvc@v2

- name: Configure DVC remote
  run: |
    dvc remote modify myremote access_key_id \
      ${{ secrets.AWS_ACCESS_KEY_ID }}
    dvc remote modify myremote secret_access_key \
      ${{ secrets.AWS_SECRET_KEY }}

- name: Pull data
  run: dvc pull

- name: Reproduce pipeline
  run: dvc repro

- name: Compare metrics
  run: |
    dvc metrics diff --md >> $GITHUB_STEP_SUMMARY
\end{minted}
\end{codeblock}

\begin{remark}
Use \cli{dvc metrics diff} in pull requests to automatically show how
a code change affects model performance. The CML (Continuous Machine
Learning) tool by Iterative can post these comparisons as PR comments.
\end{remark}

\section{CI/CD Platform Comparison}
\label{sec:cicd-comparison}

\begin{center}
\begin{tabular}{@{}lcccc@{}}
\toprule
\textbf{Criterion} & \textbf{GitHub Actions} & \textbf{GitLab CI} &
  \textbf{Jenkins} & \textbf{CML} \\
\midrule
Hosted & Yes & Yes & Self-hosted & Cloud \\
GPU runners & Paid & Yes & Self-hosted & Yes \\
Config format & YAML & YAML & Groovy/YAML & YAML \\
ML-specific & No & No & No & Yes \\
Marketplace & Large & Medium & Large & Small \\
Free tier & Generous & Generous & Free & Free \\
Best for & Open source & GitLab users & Enterprise & ML teams \\
\bottomrule
\end{tabular}
\end{center}

\section{Best Practices and Anti-patterns}

\begin{bestpractice}[title=\textbf{CI/CD best practices for ML}]
\begin{enumerate}
  \item \textbf{Fast CI}: use a small data subset for CI training
    (full training offline or nightly).
  \item \textbf{Model performance gates}: fail the pipeline if
    metrics drop below thresholds.
  \item \textbf{Cache dependencies}: use pip/conda caching to speed
    up builds.
  \item \textbf{Secrets management}: never hardcode API keys; use
    repository secrets.
  \item \textbf{Artifact management}: upload trained models as CI
    artifacts for traceability.
  \item \textbf{Separate CI config}: use a lightweight
    \file{ci\_config.yaml} with fewer epochs and smaller data.
\end{enumerate}
\end{bestpractice}

\begin{antipattern}[title=\textbf{CI/CD anti-patterns}]
\begin{itemize}
  \item Training on the full dataset in CI---pipelines take hours.
  \item No model validation step---broken models reach production.
  \item Storing large data files in the Git repository.
  \item Skipping linting and type checking for ML code.
  \item No integration tests for the serving API.
  \item Manual deployment after CI passes---defeats the purpose of CD.
\end{itemize}
\end{antipattern}

\section{Mini-project: ML CI/CD Pipeline}
\label{sec:mini-project-ch9}

\begin{miniproject}[title=\textbf{Mini-project 9: CI/CD for ML}]
\begin{enumerate}
  \item Set up \pkg{pre-commit} with black, ruff, and mypy hooks.
  \item Write a GitHub Actions workflow that:
    \begin{itemize}
      \item Runs linting and unit tests
      \item Trains a model on a small subset
      \item Validates model performance against thresholds
    \end{itemize}
  \item Add a model validation script that exits with code 1 if
    metrics are below thresholds.
  \item Configure DVC to pull test data in the CI pipeline.
  \item Add a deployment step to staging (Docker build + push to a
    registry).
  \item Verify the full pipeline runs on a pull request.
\end{enumerate}
\end{miniproject}

\section{Exercises}
\label{sec:exercises-ch9}

\begin{exercise}[$\bigstar$ --- Pre-commit setup]
Configure \pkg{pre-commit} for an ML project with hooks for:
black, ruff, mypy, trailing whitespace, and large file detection
(max 500\,KB). Run it on an existing project and fix all issues.
\end{exercise}

\begin{exercise}[$\bigstar\bigstar$ --- GitHub Actions workflow]
Write a GitHub Actions workflow for an ML project that:
\begin{enumerate}
  \item Runs on push to \texttt{main} and on pull requests
  \item Uses a matrix strategy to test on Python 3.10 and 3.11
  \item Caches pip dependencies
  \item Runs unit tests and generates a coverage report
  \item Uploads the coverage report as an artifact
\end{enumerate}
\end{exercise}

\begin{exercise}[$\bigstar\bigstar\bigstar$ --- Complete ML CI/CD]
Build an end-to-end CI/CD pipeline that:
\begin{enumerate}
  \item Pulls versioned data with DVC
  \item Trains a model and logs metrics to MLflow
  \item Compares metrics against the current production model
  \item Builds a Docker image and pushes to GitHub Container Registry
  \item Deploys to a staging environment
  \item Posts a metric comparison table as a PR comment using CML
\end{enumerate}
\end{exercise}

\section{Cheatsheet}

\begin{cheatsheet}[title=\textbf{Chapter 9 Summary}]
\begin{tabular}{@{}p{4cm}p{8.5cm}@{}}
\toprule
\textbf{Concept} & \textbf{Key Point} \\
\midrule
CI for ML & Lint + test + train (small) + validate metrics \\
CD for ML & Auto-deploy if model passes validation gates \\
Pre-commit & Local hooks: black, ruff, mypy, large file check \\
GitHub Actions & YAML workflows in \file{.github/workflows/} \\
Model validation & Script that fails CI if accuracy/F1 too low \\
DVC in CI & \cli{dvc pull} + \cli{dvc repro} + \cli{dvc metrics diff} \\
Secrets & Use repository secrets, never hardcode \\
\bottomrule
\end{tabular}
\end{cheatsheet}
